\subsection{Datasets}

We evaluate our method on two tool-use benchmarks: \textbf{ToolHop}\footnote{\url{https://hf.co/datasets/bytedance-research/ToolHop}, CC BY 4.0 license}~\cite{ye2025toolhop} and
\textbf{BFCL}\footnote{\url{https://hf.co/datasets/gorilla-llm/Berkeley-Function-Calling-Leaderboard}, Apache license 2.0}~\cite{berkeley-function-calling-leaderboard}. ToolHop consists of compositional,
multi-turn queries requiring multi-step tool reasoning, while the \texttt{live-multiple}
subset of BFCL provides a complementary single-step setting with a larger and more diverse
tool set, making tool selection more challenging.

For both datasets, we assume access to ground-truth execution traces during training, but no
ground-truth rules are provided. To evaluate generalization, we define two splits:
\textbf{\rand}, a random in-distribution split, and \textbf{\unseen}, a held-out split
containing queries with tools not seen during training. Dataset statistics are summarized
in Table~\ref{tab:dataset}.



